To thoroughly evaluate our proposed framework, we conduct both simulation and real-world experiments with three objectives:
(i) to evaluate the effectiveness and robustness of the proposed multimodal CFM policy;
(ii) to demonstrate the real-time inference capability of the distilled one-step student policy in dynamic scenarios; and
(iii) to verify that the student faithfully captures the multimodal action distributions learned by the teacher.

\subsection{Simulation Experiment}

We conduct simulation experiments on the RLBench benchmark~\cite{james2019rlbench}.
For a fair comparison with prior flow-based methods, we adopt the identical set of eight manipulation tasks used in PointFlowMatch~\cite{chisari2024learning} (Fig.~\ref{fig:RLBench}).


\begin{figure}
    \centering
    \includegraphics[width=\linewidth]{figures/RLBench.pdf}
\caption{
Eight representative manipulation tasks from the RLBench.
Tasks are ordered from left to right and top to bottom:
\textit{close door}, \textit{open box}, \textit{open fridge}, \textit{open oven},
\textit{take frame off hanger}, \textit{unplug charger},
\textit{take shoes out of box}, and \textit{put books on bookshelf}.
}
    \label{fig:RLBench}
\end{figure}


\subsubsection{Implementation Details} 
\paragraph{Baselines.}
We compare our method against representative imitation learning approaches, including: (i) diffusion-based policies, including Diffusion Policy \cite{chi2023diffusionpolicy} and Diffusion Policy 3D \cite{Ze2024DP3};
and (ii) flow-matching-based approaches, including PointFlowMatch \cite{chisari2024learning} and FlowPolicy \cite{zhang2024flowpolicy}.


\paragraph{Data Collection.} For each task, we collect 100 expert demonstrations using the RLBench built-in motion planner. The observation space consists of multi-view RGB images (5 cameras) and aligned depth maps. Point clouds are back-projected from the depth maps and downsampled via voxelization to 4,096 points.
\paragraph{Training.}
The multimodal CFM teacher is trained for 2,000 epochs on the expert dataset, and the best checkpoint is selected based on validation success rate.
Using the frozen teacher, we generate 16 trajectory candidates of length 32 per observation to construct the distilled dataset.
The student policy is then trained on this dataset for 1,500 epochs.

\paragraph{Evaluation.} To ensure statistical significance, all models are evaluated on 300 unseen episodes with randomized object poses. All training and inference experiments are conducted on a workstation equipped with a NVIDIA RTX PRO 6000.



\subsubsection{Results}
\input{table/sim}
Quantitative Results. Table ~\ref{tab:sim} summarizes the success rates (SR) and inference speeds (IS) of our distilled student model compared to several state-of-the-art (SOTA) diffusion-based and flow-based policies.

Superiority of Single-Step Generation. A key observation from Table [X] is that while original diffusion-based methods like Diffusion Policy and DP3 suffer from a catastrophic performance drop when limited to 1-step inference (e.g., Diffusion Policy’s SR drops from 18.7\% to 0.0\%), our distilled Student model maintains a high average success rate of 67.3\% in a single step. This performance is not only significantly higher than other 1-step variants (e.g., 37.2\% for PointFlowMatch and 42.4\% for FlowPolicy) but also remains competitive with the 50-step Teacher model (72.6\% SR), retaining approximately 92.7\% of the teacher's capability.

Balance Between Precision and Speed. In terms of inference efficiency, the Student model achieves a remarkable speed of 125.00 Hz, which is a 15.4$\times$ speedup over the CFM-Teacher (8.09 Hz). Compared to other 1-step baselines, although some methods like FlowPolicy (236.43 Hz) show higher raw frequency, their success rates are substantially lower (42.4\% vs. 67.3\%). Our method strikes the best balance between real-time responsiveness and task execution precision, making it highly suitable for high-frequency closed-loop robotic control.

Task-Specific Performance. Across the eight manipulation tasks, the Student model demonstrates consistent robustness. In relatively simpler tasks like open box, the Student model achieves a 98.0\% success rate, nearly identical to the Teacher. In more complex, long-horizon tasks such as open oven and shoes box, our model significantly outperforms the 1-step FlowPolicy by 28.4\% and 30.9\% respectively. This indicates that our distillation process effectively preserves the multi-modal distribution of the action space, which is often lost in naive 1-step approximations.

Consistency across Diverse Scenarios. The results also highlight the limitations of traditional Diffusion Policy in tasks requiring high precision (e.g., open oven or books shelf where SR is near 0\%). In contrast, by leveraging Conditional Flow Matching (CFM) and our distillation strategy, our model provides a much stronger foundation for learning complex manipulation, consistently outperforming the DP3 baseline in both 50-step and 1-step configurations across almost all tasks.

\section{Ablation Studies}
\label{sec:ablation}
To investigate the contribution of individual components in our framework, we conduct ablation studies focusing on three key aspects: (1) the necessity of the trimodal sensor fusion strategy, (2) the number of distilled trajectory candidates $K$, and (3) the prediction horizon $T$. Results are averaged over the most challenging tasks and reported in Table~\ref{tab:ablation_unified}.

\begin{table}[t]
    \centering
    \small
    \setlength{\tabcolsep}{12pt} 
    \renewcommand{\arraystretch}{1.15} 
    \caption{\textbf{Ablation Studies.} We analyze the impact of sensor modalities, distillation candidates $K$, and prediction horizon $T$. Default setting: \textbf{Ours} ($K{=}16, T{=}32$, Trimodal).}
    \label{tab:ablation_unified}
    
    \begin{tabular}{lc}
        \toprule
        \textbf{Variant} & \textbf{SR (\%)} \\
        \midrule
        \multicolumn{2}{l}{\textit{(a) Effect of Modality}} \\
        \midrule
        RGB + Depth (Concatnate) & 30.4 \\
        RGB + Depth (Attention) & 39.6 \\
        RGB + Depth (Gated Attention) & 66.2 \\
        PCD + RGB & 60.7 \\
        PCD + Depth & 59.8 \\
        \rowcolor{gray!10} \textbf{Ours (Trimodal)} & \textbf{74.1} \\ 
        
        \midrule
        \multicolumn{2}{l}{\textit{(b) Effect of Distillation Candidates $K$}} \\
        \midrule
        $K = 1$ & 47.3 \\
        $K = 4$ & 58.9 \\
        $K = 10$ & 68.1 \\
        \textbf{$K = 16$ (Ours)} & \textbf{74.1} \\
        
        \midrule
        \multicolumn{2}{l}{\textit{(c) Effect of Prediction Horizon $T$}} \\
        \midrule
        $T = 1$ & 30.7 \\
        $T = 8$ & 64.8 \\
        $T = 16$ & 70.9 \\
        \textbf{$T = 32$ (Ours)} & \textbf{74.1} \\
        \bottomrule
    \end{tabular}
\end{table}

\paragraph{Impact of Multi-Modal Fusion.}
We first evaluate the effectiveness of our trimodal perception design, as shown in Table~\ref{tab:ablation_unified}(a).
The results clearly demonstrate the superiority of combining Point Cloud (PCD), RGB, and Depth.
Relying solely on 2D modalities (RGB + Depth) yields suboptimal performance, even when using advanced fusion mechanisms like attention ($39.6\%$) or gating ($66.2\%$). This suggests that explicit 3D geometric information provided by point clouds is crucial for precise manipulation.
Furthermore, comparing dual-modal baselines (PCD+RGB at $60.7\%$ and PCD+Depth at $59.8\%$) with our full trimodal approach (\textbf{72.6\%}) confirms that neither visual texture nor depth geometry alone is sufficient; their complementary integration achieves the best robustness.
Notably, our \textit{Attention+Gate} fusion mechanism significantly outperforms naive concatenation ($30.4\%$) and standard attention ($39.6\%$), validating the effectiveness of our specific architecture in selecting informative features.

\paragraph{Effect of Distillation Candidates $K$.}
In Table~\ref{tab:ablation_unified}(b), we analyze the sensitivity of the model to the number of distilled trajectory candidates $K$.
We observe a monotonic improvement in Success Rate (SR) as $K$ increases from 1 to 16.
Using a single deterministic candidate ($K=1$) results in a low SR of $47.3\%$, likely because the student policy fails to fully capture the multi-modal distribution of the teacher's behavior.
Increasing $K$ allows the model to approximate the teacher's action distribution more accurately by covering diverse modes.
The performance improves to $68.1\%$ at $K=10$ and reaches its peak at $K=16$ with \textbf{72.6\%}, indicating that $16$ candidates are sufficient for high-fidelity behavior cloning via distillation without incurring excessive computational overhead.

\paragraph{Importance of Prediction Horizon $T$.}
Finally, Table~\ref{tab:ablation_unified}(c) investigates the impact of the action prediction horizon $T$.
The results highlight the importance of temporal consistency in manipulation tasks.
A greedy, single-step prediction ($T=1$) leads to poor performance ($30.7\%$), as it lacks foresight and temporal smoothness, often resulting in jerky robot motions.
As we extend the horizon to $T=8$ and $T=16$, the success rate improves significantly to $64.8\%$ and $70.9\%$, respectively.
The best performance is achieved at $T=32$ (\textbf{72.6\%}), confirming that predicting longer-term trajectories enables the agent to plan more coherent motions and effectively handle complex multi-stage tasks.

\subsection{Real-World Experiments}
We evaluate the generalization capability and real-time performance of our method in real-world robotic manipulation scenarios. To this end, we deploy our policy on a physical robot platform equipped with a multi-view RGB-D perception system and design a set of challenging manipulation tasks involving dynamic objects, multi-stage execution, and human-induced disturbances.

Specifically, we consider five representative real-world tasks that probe different aspects of multimodal perception, trajectory diversity, and reactive control:
\begin{enumerate}
    \item \textbf{Dynamic Cube Stowing:} 
    The robot grasps one cube from three adjacent cubes while the target cube is perturbed twice by human interference. It then places the cube into a box, which is also perturbed once during the placement stage.

    \item \textbf{Kitchen Microwave Loading:} 
    The robot picks up food from the stovetop, places it into a container, grasps the container, and transfers it into a microwave oven.

    \item \textbf{Kitchen Cleanup:} 
    The robot grasps a pot from the stovetop and places it into the sink, then grasps a soap bottle from the shelf and places it into the sink.

    \item \textbf{Dynamic Cabinet Opening:} 
    The robot approaches a vertically moving cabinet door, grasps the handle at an appropriate moment, and pulls the door open while compensating for the door’s dynamic motion.

    \item \textbf{Dynamic Grasping:} 
    The robot reaches into a moving cabinet, grasps a soap bottle, and retrieves it while adapting to the cabinet’s motion.
\end{enumerate}

Figure~\ref{fig:real_setup} illustrates the real-world experimental setup and task groupings. 
Task~(a) corresponds to \emph{Dynamic Cube Stowing}. 
Task~(b) includes two kitchen manipulation scenarios: \emph{Microwave Loading} and \emph{Kitchen Cleanup}. 
Task~(c) includes two cabinet-related tasks: \emph{Dynamic Cabinet Opening} and \emph{Dynamic Grasping}.

\begin{figure*}[t]
    \centering
    \begin{subfigure}[t]{0.32\linewidth}
        \centering
        \includegraphics[width=\linewidth]{figures/Setup_Cube.pdf}
        \caption{Dynamic Cube Stowing}
        \label{fig:sub1}
    \end{subfigure}
    \hfill
    \begin{subfigure}[t]{0.32\linewidth}
        \centering
        \includegraphics[width=\linewidth]{figures/Setup_Kitchen.pdf}
        \caption{Kitchen Tasks}
        \label{fig:sub2}
    \end{subfigure}
    \hfill
    \begin{subfigure}[t]{0.32\linewidth}
        \centering
        \includegraphics[width=\linewidth]{figures/Setup_Cabinet.pdf}
        \caption{Cabinet Tasks}
        \label{fig:sub3}
    \end{subfigure}
    \caption{
    Real-world experimental setup. The robot is equipped with a multi-view RGB-D perception system, consisting of a wrist-mounted camera and multiple external cameras, which provides comprehensive visual coverage for dynamic manipulation. Target objects and task-relevant regions are highlighted for clarity.
    }
    \label{fig:real_setup}
\end{figure*}


\subsubsection{Implementation Details}

\paragraph{Hardware Setup.}
All real-world experiments are conducted on a 7-DoF Franka Emika Panda robotic arm equipped with a parallel gripper. 
The perception system consists of a wrist-mounted RGB-D camera and multiple external RGB-D cameras, providing multi-view observations of the workspace.

\paragraph{Task Execution and Data Collection.}
We collect expert demonstrations for each real-world task via teleoperation using a 3D SpaceMouse.
For each task, 300 demonstrations are recorded, exhibiting substantial trajectory diversity arising from variations in object initial configurations, grasping poses, and motion paths.

Importantly, each task can be naturally decomposed into a sequence of semantic sub-stages (e.g., reaching, grasping, placing), whose number and ordering may differ across tasks.
Figure~\ref{fig:traj_single_column} visualizes representative execution trajectories for all five tasks, illustrating that even within the same task, multiple valid trajectories and stage transitions can occur.
This inherent variability motivates modeling a multi-modal trajectory distribution rather than a single deterministic policy.

\begin{figure}[t]
    \centering
    \begin{subfigure}{\columnwidth}
        \centering
        \includegraphics[width=\linewidth, height=0.16\textheight, keepaspectratio]{figures/Traj_Cube.pdf}
        \caption{Dynamic Cube Stowing}
    \end{subfigure}
    \vspace{0.4em}
    \begin{subfigure}{\columnwidth}
        \centering
        \includegraphics[width=\linewidth, height=0.16\textheight, keepaspectratio]{figures/Traj_Kitchen_Banana.pdf}
        \caption{Kitchen Microwave Loading}
    \end{subfigure}
    \vspace{0.4em}
    \begin{subfigure}{\columnwidth}
        \centering
        \includegraphics[width=\linewidth, height=0.16\textheight, keepaspectratio]{figures/Traj_Kitchen_Soap.pdf}
        \caption{Kitchen Cleanup}
    \end{subfigure}
    \vspace{0.4em}
    \begin{subfigure}{\columnwidth}
        \centering
        \includegraphics[width=\linewidth, height=0.16\textheight, keepaspectratio]{figures/Traj_Cabinet_Open_door.pdf}
        \caption{Dynamic Cabinet Opening}
    \end{subfigure}
    \vspace{0.4em}
    \begin{subfigure}{\columnwidth}
        \centering
        \includegraphics[width=\linewidth, height=0.16\textheight, keepaspectratio]{figures/Traj_Cabinet_Grasp_Soap.pdf}
        \caption{Dynamic Grasping}
    \end{subfigure}
    \caption{
    Visualization of trajectory execution and diversity for five real-world manipulation tasks.
    Each subfigure overlays multiple expert demonstrations, illustrating diverse motion patterns
    and task-dependent multi-stage behaviors under varying object configurations.
    }
    \label{fig:traj_single_column}
\end{figure}
During data collection, multi-view RGB images and aligned depth maps are captured.
A point cloud is generated from a single depth view, cropped to the tabletop workspace, and downsampled to 4096 points using voxel grid filtering.


\paragraph{Training Procedure.}
The training pipeline follows the same procedure as in simulation.
We first train the multimodal Conditional Flow Matching (CFM) teacher on the collected real-world demonstrations.
The teacher is then used to generate a distilled multi-hypothesis trajectory dataset via flow-based sampling.
Finally, a one-step student policy is trained using IMLE with a Chamfer-based diversity loss.
All network architectures and hyperparameters are kept identical to those used in simulation experiments.

\paragraph{Deployment.}
Policy inference is performed on a workstation equipped with an NVIDIA RTX 3070 GPU.
Predicted actions are transmitted via Ethernet to a dedicated control PC running a real-time kernel, which executes high-frequency control commands on the robot.
For the dynamic cabinet opening task, an additional robotic arm induces continuous vertical motion of the cabinet door, controlled by a separate laptop with an NVIDIA RTX 4080 GPU.

\paragraph{Evaluation Protocol.}
For each task, we conduct 30 real-world evaluation trials and report the average task success rate together with the policy inference frequency (Hz).




\subsubsection{Results Analysis}
\input{table/real}
Overall Performance and Real-Time Capability.As shown in Table~\ref{tab:comparison}, our distilled \textbf{Student} policy achieves a compelling balance between task success rate and inference speed in real-world scenarios. With an average inference speed (\textbf{Avg IS}) of \textbf{125.0 Hz}, our method is approximately \textbf{43$\times$ faster} than the CFM-Teacher (2.9 Hz) and 35$\times$ faster than the PointFlowMatch 50-step baseline (3.6 Hz). This high-frequency control capability is crucial for the dynamic tasks we designed, such as \textit{Dynamic Cabinet Opening} and \textit{Dynamic Grasping}, where the robot must react instantaneously to moving targets and disturbances.Comparison with Baselines.A critical finding is the comparison with the \textit{PointFlowMatch (1-step)} baseline. While the 1-step baseline achieves a high inference speed (210.0 Hz), its performance collapses completely, yielding a near-zero average success rate (\textbf{2.7\%}). This indicates that naive one-step generation fails to capture the complex multi-modal distributions required for these manipulation tasks. In stark contrast, our \textbf{Student} model maintains a robust average success rate of \textbf{71.3\%} using single-step inference. This demonstrates that our distillation method effectively compresses the knowledge of the teacher into a fast-reactive policy without the catastrophic performance drop seen in other methods.Preservation of Teacher Proficiency.Although the CFM-Teacher (50-step) achieves the highest average success rate (84.0\%), its slow inference speed (2.9 Hz) creates significant latency, causing "stop-and-go" behavior that is undesirable in dynamic human-robot interaction settings. Our Student model retains \textbf{$\sim$85\%} of the Teacher's performance (71.3\% vs. 84.0\%) while unlocking real-time responsiveness. Notably, in long-horizon tasks like \textit{Kitchen Cleanup}, the Student achieves an 86.7\% success rate, closely matching the Teacher's 96.7\%, proving its ability to handle multi-stage execution effectively.Robustness in Dynamic Scenarios.The results on dynamic tasks highlight the practical benefits of our method.\begin{itemize}\item In \textbf{Dynamic Cube Stowing}, despite heavy human interference (two perturbations during grasping and one during placement), the Student achieves a \textbf{76.7\%} success rate. The 125 Hz control loop allows the robot to rapidly re-plan and track the cube immediately after it is pushed away, a behavior that is difficult for the 2.9 Hz Teacher to execute smoothly despite its higher theoretical success rate.\item In \textbf{Dynamic Cabinet Opening} and \textbf{Dynamic Grasping}, where the environment itself is in motion, the Student achieves 43.3\% and 66.7\% success rates, respectively. While there is a performance gap compared to the 50-step Teacher (60.0\% and 80.0\%), the Student remains the only viable 1-step solution, as the baseline 1-step method fails completely (0.0\%) in these precision-demanding dynamic tasks.\end{itemize}